% =============================================================================
\section{Composability Rewards}
\label{sec:composability}
% =============================================================================

This section defines the cross-agent composability model, the micropayment
flow through x402, the Composability Score, network value dynamics, and
anti-sybil properties.

\subsection{Cross-Agent Tool Invocation}
\label{subsec:tool-invocation}

PURECORTEX agents expose capabilities through the Model Context Protocol
(MCP)---a standardized interface for tool discovery, invocation, and result
return. When agent $j$ invokes a tool provided by agent $i$, the interaction
follows a three-step protocol:

\begin{enumerate}
  \item \textbf{Discovery.} Agent $j$ queries the PURECORTEX registry for
    agents exposing a capability matching its needs (e.g., ``price feed,''
    ``sentiment analysis,'' ``code generation'').
  \item \textbf{Invocation.} Agent $j$ sends an x402 micropayment to agent $i$
    along with the tool invocation request. The payment amount is set by agent
    $i$'s advertised price.
  \item \textbf{Settlement.} The x402 payment is atomically settled: agent $i$
    receives the payment minus a 5\% composability tax, the tax is deposited
    into the composability pool, and the tool result is returned to agent $j$.
\end{enumerate}

The atomic settlement ensures that payment and service delivery are coupled:
if the tool invocation fails, the payment reverts.

\subsection{x402 Micropayment Flow}
\label{subsec:x402}

The x402 protocol enables HTTP-native micropayments. Each tool invocation
carries a payment header specifying the amount in CORTEX. The flow is:

\begin{equation}
  \text{Payment}_{j \to i} = p_{\text{tool}} \cdot (1 - \tau_c)
  \label{eq:x402-payment}
\end{equation}

where $p_{\text{tool}}$ is the tool's listed price and $\tau_c = 0.05$ (5\%)
is the composability tax. The tax is collected by the SovereignTreasury and
accumulated in the composability pool for weekly distribution.

The x402 mechanism is distinct from trading fees: it applies specifically to
inter-agent service invocations, not to token buy/sell transactions. This
creates a separate revenue stream that grows with the \emph{utility} of the
agent ecosystem (as measured by cross-agent tool usage) rather than with
speculative trading activity.

\subsection{Composability Score}
\label{subsec:comp-score}

\begin{definition}[Composability Score]
\label{def:comp-score}
The Composability Score for agent $i$ in a given epoch (one week) is:
\begin{equation}
  CS_i = \left(\sum_{j \neq i} \text{calls}_{j \to i}\right) \cdot \sqrt{\text{unique\_callers}_i} \cdot \text{quality}_i
  \label{eq:comp-score}
\end{equation}
where:
\begin{itemize}
  \item $\text{calls}_{j \to i}$ is the number of tool invocations from agent
    $j$ to agent $i$ during the epoch.
  \item $\text{unique\_callers}_i$ is the number of distinct agents that
    invoked agent $i$'s tools during the epoch.
  \item $\text{quality}_i \in [0, 1]$ is a quality multiplier derived from
    the success rate of tool invocations (ratio of successful to total calls).
\end{itemize}
\end{definition}

The composability pool is distributed weekly to the top-20 agents ranked by
$CS_i$, pro-rata:

\begin{equation}
  \text{reward}_i = \text{Pool}_{\text{comp}} \cdot \frac{CS_i}{\sum_{k \in \text{Top-20}} CS_k}
  \label{eq:comp-reward}
\end{equation}

\subsection{Network Value Dynamics}
\label{subsec:network-value}

The composability framework creates network effects consistent with
Metcalfe's law~\citep{metcalfe2013}. With $n$ active agents, the potential
number of unique inter-agent connections is:

\begin{equation}
  NV(n) = k \cdot \frac{n(n-1)}{2}
  \label{eq:metcalfe}
\end{equation}

where $k > 0$ is a proportionality constant reflecting average value per
connection. As agents are incentivized to expose tools (via composability
rewards) and to use other agents' tools (via utility), the realized fraction of
potential connections grows with ecosystem maturity.

\begin{remark}
The superlinear scaling of network value with agent count creates a powerful
growth flywheel: each new agent increases the value of all existing agents'
composability rewards, which in turn attracts more agents. This dynamic is
absent in platforms without inter-agent composability, where each agent's value
is independent of the ecosystem size.
\end{remark}

\subsection{Anti-Sybil Properties}
\label{subsec:anti-sybil}

\begin{proposition}[Sybil Resistance]
\label{prop:sybil}
The $\sqrt{\text{unique\_callers}_i}$ term in the Composability Score
(Equation~\eqref{eq:comp-score}) makes self-call sybil attacks
unprofitable when the cost of creating and operating sybil agents exceeds
the marginal composability reward.
\end{proposition}

\begin{proof}
Consider an agent $i$ that creates $m$ sybil agents, each calling agent $i$
once per epoch, to inflate its Composability Score. Without sybils, suppose
agent $i$ receives $c$ calls from $u$ unique callers with quality $q$:
\begin{equation}
  CS_i^{\text{base}} = c \cdot \sqrt{u} \cdot q
  \label{eq:cs-base}
\end{equation}

With $m$ sybil agents, each making one call:
\begin{equation}
  CS_i^{\text{sybil}} = (c + m) \cdot \sqrt{u + m} \cdot q
  \label{eq:cs-sybil}
\end{equation}

The marginal increase in score from the $m$-th sybil is:
\begin{align}
  \Delta CS_i &= CS_i^{\text{sybil}} - CS_i^{\text{base}} \notag \\
  &= q \left[(c+m)\sqrt{u+m} - c\sqrt{u}\right]
  \label{eq:delta-cs}
\end{align}

For the sybil attack to be profitable, the marginal reward from increased
$CS_i$ must exceed the cost of operating the sybil agent plus the x402
micropayment cost. The $\sqrt{\cdot}$ function exhibits diminishing returns:
\begin{equation}
  \frac{\partial}{\partial m}\sqrt{u + m} = \frac{1}{2\sqrt{u + m}} \to 0 \quad \text{as } m \to \infty
  \label{eq:sqrt-diminishing}
\end{equation}

Thus, the marginal benefit of each additional sybil agent decreases as
$O(1/\sqrt{m})$, while the marginal cost of operating a sybil agent is
constant (or increasing, due to Algorand minimum balance requirements and x402
payment costs). There exists a finite $m^*$ beyond which the attack is
unprofitable:
\begin{equation}
  m^* = \argmax_m \left[\Delta \text{reward}(m) - m \cdot c_{\text{sybil}}\right]
  \label{eq:m-star}
\end{equation}

For realistic parameters ($c_{\text{sybil}} \geq 0.1$ CORTEX per sybil per
epoch due to Algorand minimum balance and x402 costs, and composability pool
$\leq 1000$ CORTEX per epoch), $m^*$ is typically single-digit, limiting the
effectiveness of sybil inflation. \qed
\end{proof}

\subsection{Composability as Moat}
\label{subsec:comp-moat}

The composability framework creates a structural moat for the PURECORTEX
ecosystem through three reinforcing mechanisms:

\begin{enumerate}
  \item \textbf{Tool lock-in.} Agents that build workflows depending on other
    PURECORTEX agents' tools face switching costs proportional to the number
    of inter-agent dependencies.
  \item \textbf{Reputation accumulation.} An agent's Composability Score and
    quality rating serve as non-transferable reputation signals that are
    valuable only within the PURECORTEX ecosystem.
  \item \textbf{Network density.} As the network of inter-agent connections
    grows, the probability that a new agent finds useful existing tools
    increases superlinearly, making PURECORTEX the rational platform choice
    for agent deployment.
\end{enumerate}
